\begin{center}
\LARGE
Volume of Parameterized Polyhedra 

\end{center}

\begin{center}
\Large
Hoon Hong, Manfred Minimair

\end{center}

\noindent
Date:  July 19th (Friday)\\
Time:  16:00-16:30\\

\begin{center}
\large
Abstract
\end{center}

In this talk, we study the problem of computing the volume of the polyhedron defined by a given
system of linear equalities. In fact, we study a parametric version of the problem: given a
parametric system of linear inequalities, to find an expression (in the parameters) for the volume. 

This problem naturally arise in automatic parallelization of programs (in particular vectorization
of loops). It is also needed in computer aided design, where some mechanical parts are modeled by
polyhedron and we need to know their volumes for cost/weight computation. 

We describe an algorithm (improvement/completion of the previous algorithms). It essentially
carries out recursion on the variables, where each recursion consists of projection (variables
elimination) and lifting (symbolic integration in suitable measure). In order to control the huge
intermediate result swell, inconsistent and redundant computational branches are detected and
eliminated. 

